In many practical applications, a machine learning model does not receive
all training data at once. New data arrives step by step, and a model that
is updated on new data usually forgets what it has learned before. This
problem is called catastrophic forgetting. The goal of this project is to
investigate whether a Mixture of Experts model, in which a small gating
network routes every input to a few specialised sub networks, forgets less
than a normal neural network when tasks are learned one after another.

For this purpose, a complete experimental framework was implemented in
Python and PyTorch. Two standard benchmarks were built from the MNIST data
set: SplitMNIST with five two class tasks and PermutedMNIST with five
random pixel permutations. A naive multi layer perceptron, the same network
with Elastic Weight Consolidation, a wider network with more parameters,
and the proposed Mixture of Experts model were trained sequentially and
compared with the metrics average accuracy, average forgetting and backward
transfer. Every experiment was repeated with three random seeds.

The results show that the Mixture of Experts model forgets clearly less
than the naive baseline on SplitMNIST. On PermutedMNIST, the first version
of the model performed worse than expected. A detailed analysis showed two
reasons: the load balancing loss forces all tasks to share all experts, and
more importantly, the task specific information of PermutedMNIST is stored
in the shared input layer, which the router cannot protect. A second
architecture with routing directly at the input solves this problem and
reaches the best knowledge retention of all compared approaches on both
benchmarks. An analysis of the gating network confirms that different tasks
are routed to different experts, which explains the reduced forgetting.
